\subsection{Timing pilot and the full-cohort decision}
\label{sec:compute_details}
The label-free timing pilot selected two proteins by seeded hash from each tercile of maximum 13~\AA\ support size, plus the preset principal case 1ULR. The selected identifiers were 1MF7, 1PMY, 1ULR, 2NUH, 2WUJ, 3M3P, and 3T47. The pilot evaluated the complete feature suite, including graph/native descriptors, degree-zero and degree-one operators, eigensolvers, and serialization. Its recorded wall time was 37.96 seconds. The pilot preceded restoration of modified peptide coordinates: extrapolating its slowest per-residue cost to the initial 78,223-residue audit, allowing four workers, and multiplying by two yielded a projected full-suite time of 1,773.57 seconds. This was below the declared 21,600-second degree-one feature budget by more than an order of magnitude. The decision selected the full 364-protein eligible cohort, whose final residue count is 78,400. The retained pilot values describe a planning projection; actual production runtime is reported separately.

Nested 60- and 20-protein fallback cohorts were prepared without using prediction outcomes, but were not selected by the pilot decision. Within every reported comparison, models use the same completed cohort and frozen fold assignments. Failed or uncompleted parameter combinations are not represented by zero features or zero performance. Runtime reports distinguish feature generation from fitting and explanation. Timing measurements depend on the stated implementation and hardware and do not establish an architecture-independent complexity bound.
